%% ============================================================
%% CHAPTER 3 — MimiClaw: an Agent Runtime Without an OS
%%                                        (budget: ~9,000 words)
%% Role: the system chapter. Source material: docs/mimiclaw-system-report.md
%% (expand each report section ~2x, add figures).
%% ============================================================
\chapter{MimiClaw: an Agent Runtime Without an Operating System}
\label{ch:system}

\section{Design Goals and Non-Goals}
\label{sec:sys-goals}
% Goals: complete runtime semantics (loop/tools/memory/channels/autonomy),
% survives unattended operation, safety boundary in deterministic code,
% fits alongside WiFi+TLS in 512KB SRAM. Non-goals: on-device inference,
% multi-agent, hard real-time guarantees.

\section{The Five Constraints of Bare Metal}
\label{sec:sys-constraints}
% One subsection per constraint (no processes; tiered scarce memory; flat
% filesystem; no dynamic runtime; failure as normal case). Each ends with
% a forward pointer to the mechanism that answers it.

\section{Architecture: a Message Bus on a Dual-Core Split}
\label{sec:sys-architecture}
% Figure: block diagram (channels -> inbound queue -> agent loop (core 1)
% -> outbound queue -> dispatch (core 0)). Ownership-transfer discipline;
% backpressure = drop-and-log. The asyncio-to-FreeRTOS translation table.

\section{Memory Tiering and Degraded Startup}
\label{sec:sys-memory}
% PSRAM-vs-SRAM placement rule; fixed lifetime buffers (16/32/8 KB);
% compile-time footprint bound. The descending stack ladder as a worked
% example of graceful degradation.

\section{Persistent Memory: the File Paradigm on Flat Flash}
\label{sec:sys-filememory}
% SOUL/USER/MEMORY/daily/sessions layout on SPIFFS; the model as the writer
% of its own memory; the fixed-budget context assembly (16KB, itemized
% caps); session slimming; progressive disclosure of skills.

\section{Tools and Deterministic Guards}
\label{sec:sys-guards}
% Tool registry; gpio_policy allowlist; parameter patching with the cron
% case study. State the generalizable pattern: any field the firmware
% knows with certainty is decided by the firmware.

\section{Autonomy and Fault Recovery}
\label{sec:sys-autonomy}
% cron + heartbeat; the recover-by-rebuild ladder (NVS erase, SPIFFS
% reformat, captive-portal fallback); OTA. What survives a hard reset.

\section{Resource Ledger}
\label{sec:sys-ledger}
% Tables: flash partition map; SRAM/PSRAM static budget; context budget.
% The headline number: the architectural part of an agent in <60KB SRAM.

\section{Design Principles}
\label{sec:sys-principles}
% The four principles (behavior is data; every resource has a budget; the
% firmware does not trust the model; degradation over perfection), each
% with its instances cross-referenced to earlier sections.

\section{Summary}
\label{sec:sys-summary}
% Answers the qualitative half of RQ1; the quantitative half is Ch5.
